% ch1.tex -- Module 7 ESSENTIALS, Chapter 1: The electromagnetic spectrum
% (Maxwell, the speed of light and spectroscopy). Dot points M7.IQ1.01-M7.IQ1.06.
% Knowledge points only (no questions). Condensed from src/mod7/ch1.tex:
% same notation, same facts. Figures: figures/ch1_*.pdf (copied from the full book)
% and TikZ copied from the full chapter.
\chapter[The Electromagnetic Spectrum]{The Electromagnetic Spectrum: Maxwell, the Speed of Light and Spectroscopy}\label{ch1:chap}
\covers{M7.IQ1.01, M7.IQ1.02, M7.IQ1.03, M7.IQ1.04, M7.IQ1.05, M7.IQ1.06}

\begin{summarybox}{In this chapter}
\begin{itemize}
  \item Maxwell's unification; how EM waves are made and how they travel.
  \item How the speed of light was measured, and how it now defines the metre.
  \item How spectra identify elements and tell us about the stars.
\end{itemize}
\end{summarybox}

%% ==================================================================
\section{Maxwell and electromagnetic waves}\label{ch1:sec:maxwell}
\covers{M7.IQ1.01, M7.IQ1.02}

\begin{bigidea}
Maxwell unified electricity and magnetism. His theory predicted electromagnetic (EM) waves
that travel at the speed of light. So light is an EM wave.
\end{bigidea}

\begin{keywords}
  \kw{Unification}{showing that two things that look different are parts of one thing.}
  \kw{Electromagnetic (EM) wave}{a travelling wave of electric and magnetic fields that change
    together.}
  \kw{Accelerating charge}{a charge that speeds up, slows down or changes direction.}
  \kw{Transverse wave}{a wave whose vibrations are at \ang{90} to its direction of travel.}
  \kw{In phase}{reaching their peaks at the same time and the same place.}
\end{keywords}

\textbf{Before the 1860s}, some links were known (Oersted, 1820; Faraday, 1831), but there was
no single theory. Light was thought to be a wave in the \term{luminiferous aether}, an invisible
material filling all space. From 1861 to 1865, Maxwell wrote all of electricity and magnetism as
one set of laws, now written as four laws:
\begin{enumerate}
  \item \textbf{Gauss's law (electric):} charges make electric fields (lines go from $+$ to $-$).
  \item \textbf{Gauss's law for magnetism:} field lines form closed loops; no single N or S pole.
  \item \textbf{Faraday's law:} a changing magnetic field makes an electric field, even in empty
        space.
  \item \textbf{Ampère--Maxwell law:} a current \emph{and} a changing electric field make a
        magnetic field.
\end{enumerate}

\begin{itemize}
  \item \textbf{Maxwell's new idea} was the second half of law 4. Now each changing field makes
        the other: one \term{electromagnetic field}. This is the unification of
        electricity and magnetism.
  \item \textbf{Prediction 1: EM waves.} The changing fields keep making each other and travel
        away, with no wire and no charges, at any frequency (most kinds are invisible).
  \item \textbf{Prediction 2: their speed.} It comes from the \term{permeability of free space}
        ($\mu_0$, for magnetic fields in a vacuum) and the \term{permittivity of free space}
        ($\varepsilon_0$, for electric fields). Both are measured with currents and charges. Maxwell
        got about \qty{3.1e8}{m.s^{-1}} (1862); Fizeau had measured light at about
        \qty{3.1e8}{m.s^{-1}} (1849). So \textbf{light is an EM wave}.
  \item Maxwell only \emph{predicted} (and still believed in the aether). \textbf{Hertz (1887)}
        made and detected radio waves with sparks. They reflected, refracted and were polarised,
        like light, and travelled at about $c$: this confirmed Maxwell's theory.
\end{itemize}

\begin{formulacard}[notonsheet]{Speed of electromagnetic waves (Maxwell)}
  \cardformula{c = \frac{1}{\sqrt{\mu_0\varepsilon_0}}}
  \begin{symbols}
    \sym{$c$}{speed of EM waves in a vacuum}{\si{m.s^{-1}}}
    \sym{$\mu_0$}{permeability of free space, $4\pi\times10^{-7}$ (data sheet)}{\si{N.A^{-2}}}
    \sym{$\varepsilon_0$}{permittivity of free space, \num{8.854e-12} (data sheet)}%
      {\si{A^2.s^4.kg^{-1}.m^{-3}}}
  \end{symbols}
  \whentouse{Maxwell's predicted speed of EM waves (\qty{2.998e8}{m.s^{-1}}).}
  \derive{laws 3 and 4 combine into a wave with $v^2\mu_0\varepsilon_0 = 1$.}
  \trap{forgetting the square root gives \num{8.99e16}, which is $c^2$, not $c$.}
\end{formulacard}

\subsection{How EM waves are made and how they travel}

\begin{itemize}
  \item No wave from a charge at rest or at a \emph{constant} velocity. An \emph{accelerating}
        charge disturbs its field; this disturbance spreads out at $c$.
  \item An oscillating charge sends out an EM wave of the \emph{same frequency}. A radio
        transmitter's AC supply pushes electrons back and forth in the \term{antenna} (a metal rod
        that sends out or picks up radio waves). In a receiving antenna, the wave makes a small
        AC current.
\end{itemize}

\begin{sidebyside}[0.49]
  \centering
  \begin{tikzpicture}[x={(0.82cm,0cm)}, y={(0cm,0.82cm)}, z={(-0.24cm,-0.30cm)}]
    \def\lam{3.0}
    \fill[hscseriesA!7] (0,-1.25,0) -- (6,-1.25,0) -- (6,1.25,0) -- (0,1.25,0) -- cycle;
    \fill[hscforce!15] (0,0,-1.25) -- (6,0,-1.25) -- (6,0,1.25) -- (0,0,1.25) -- cycle;
    \foreach \x in {0.125,0.25,...,6.01}
      \draw[hscseriesA!55, line width=0.35pt] (\x,0,0) -- (\x,{1.25*sin(360*\x/\lam)},0);
    \draw[hscseriesA, line width=1.3pt, domain=0:6, samples=121, smooth]
      plot (\x,{1.25*sin(360*\x/\lam)},0);
    \foreach \x in {0.125,0.25,...,6.01}
      \draw[hscforce!55, line width=0.35pt] (\x,0,0) -- (\x,0,{1.25*sin(360*\x/\lam)});
    \draw[hscforce, line width=1.3pt, domain=0:6, samples=121, smooth]
      plot (\x,0,{1.25*sin(360*\x/\lam)});
    \draw[axisline] (0,0,0) -- (6.9,0,0) node[below, diagramlabel, align=center]
      {direction\\of travel};
    \draw[axisline] (0,0,0) -- (0,1.7,0);
    \draw[axisline] (0,0,0) -- (0,0,2.0);
    \node[diagramlabel, text=hscseriesA, anchor=west] at (4.05,1.05,0) {$E$};
    \node[diagramlabel, text=hscforce, anchor=north east] at (0.65,0,1.3) {$B$};
    \draw[<->, line width=0.6pt, hscink] (0.75,1.62,0) -- (3.75,1.62,0)
      node[midway, above, diagramlabel] {$\lambda$};
    \draw[guide] (0.75,1.3,0) -- (0.75,1.68,0);
    \draw[guide] (3.75,1.3,0) -- (3.75,1.68,0);
    % one pair of crests: E and B are largest at the same place (in phase)
    \draw[hscink, line width=0.6pt, dashed] (3.75,1.25,0) -- (3.75,0,0) -- (3.75,0,1.25);
    \fill[hscink] (3.75,1.25,0) circle (1.6pt);
    \fill[hscink] (3.75,0,1.25) circle (1.6pt);
    \draw[hscforce, line width=0.5pt] (3.75,0,1.25) -- ++(0cm,-0.66cm)
      node[below, inner sep=1pt, diagramlabel, text=hscforce] {$B$ largest};
  \end{tikzpicture}
  \captionof{figure}{An EM wave at one instant: $E$ (blue, vertical plane), $B$ (red,
    horizontal plane, drawn coming out towards you). Black dots: both largest at the same
    place.}%
  \label{ch1:fig:emwave}
\nextside
  \begin{itemize}
    \item $E$ and $B$ have the same frequency and are in phase.
    \item $E$, $B$ and the direction of travel are all at \ang{90}: a transverse wave.
    \item No medium is needed. It travels at $c$ in a vacuum and carries energy.
    \item Maxwell's laws predict all of this (\cref{ch1:fig:emwave}), and the speed $c$.
  \end{itemize}
\end{sidebyside}

\subsection{The electromagnetic spectrum}

All EM waves travel at $c$ in a vacuum and differ only in frequency, and so in wavelength
(\cref{ch1:fig:emspectrum}). Visible light is about \qty{400}{nm} (violet) to \qty{700}{nm}
(red). Every source involves charges: accelerating in an antenna, or changing energy inside
atoms and nuclei.

\begin{figure}
  \centering
  \includegraphics{ch1_emspectrum.pdf}
  \caption{The electromagnetic spectrum, with a typical source of each band. Band edges are
    approximate.}\label{ch1:fig:emspectrum}
\end{figure}

\begin{formulacard}[combine={$v = f\lambda$ with $v = c$}]{Wave equation for light}
  \cardformula{c = f\lambda}
  \begin{symbols}
    \sym{$c$}{speed of light in a vacuum, \qty{3.00e8}{m.s^{-1}}}{\si{m.s^{-1}}}
    \sym{$f$}{frequency (from a period: $f = 1/T$, on the sheet)}{\si{Hz}}
    \sym{$\lambda$}{wavelength}{\si{m}}
  \end{symbols}
  \whentouse{any EM wave, from radio to gamma rays, in a vacuum or in air.}
  \derive{the sheet gives $v = f\lambda$ for every wave. For EM waves in a vacuum, $v = c$.}
  \trap{change \si{nm} to \si{m} ($\times10^{-9}$), and \si{MHz} or \si{GHz} to \si{Hz},
    before you substitute.}
\end{formulacard}

%% ==================================================================
\section{Measuring the speed of light}\label{ch1:sec:speed}
\covers{M7.IQ1.03}

\begin{bigidea}
Each new method measured a time or a distance better: machines, then atomic clocks, replaced
human timing. Since 1983, $c$ has an exact value, and the metre is defined from it.
\end{bigidea}

\begin{keywords}
  \kw{Direct method}{measure a distance and a travel time, then calculate $c = d/t$.}
  \kw{Indirect method}{calculate $c$ from other measured quantities, for example $c = f\lambda$.}
  \kw{Atomic clock}{an extremely accurate clock that counts oscillations of radiation from atoms.}
\end{keywords}

\Cref{ch1:tab:history} sums up the methods. The main ones follow.

\begin{table}
  \caption{Measuring the speed of light.}\label{ch1:tab:history}
  \centering\small
  \begin{tabular}{@{}L{24mm}L{39mm}L{31mm}L{42mm}@{}}
    \toprule
    \colhead{Who, when} & \colhead{Method} & \colhead{Result (\si{m.s^{-1}})} &
      \colhead{What limited it}\\
    \midrule
    Galileo, about 1638 & lanterns on two hills about \qty{1.5}{km} apart & no delay seen &
      reaction time is far longer than the travel time\\
    Rømer, 1676 & delays of the eclipses of Jupiter's moon Io & not instant; Huygens (1690):
      about \num{2.2e8} & eclipse timing; size of Earth's orbit\\
    Fizeau, 1849 & toothed wheel; mirror \qty{8.6}{km} away & about \num{3.1e8} & judging
      ``darkness'' by eye\\
    Foucault, 1862 & rotating mirror & \num{2.98e8} & measuring a tiny angle\\
    Michelson, 1926 & rotating 8-sided mirror; \qty{35}{km} each way & \num{2.99796e8} &
      air slows light a little; uncertainty \qty{4}{km.s^{-1}}\\
    Lasers, 1972 & measure $f$ and $\lambda$ of a laser; $c = f\lambda$ & \num{299792456}
      ($\pm1$) & the metre itself (then defined by krypton light)\\
    \bottomrule
  \end{tabular}
\end{table}

\begin{itemize}
  \item \textbf{Direct method} \hscbadgenotonsheet\ $c = d/t$. A reflected signal goes there
        \emph{and back}: $d = 2D$. Do not use $s$ for speed: on the sheet, $s$ is
        displacement.
  \item \textbf{Rømer (1676):} Io's \term{eclipses} (Io hidden in Jupiter's shadow) happen
        every 1.77 days. When Earth was on the far side of its orbit, they came up to
        \qty{22}{min} late (modern value \qty{16.6}{min}). Light takes time to cross the orbit
        ($d = 2r$).
  \item \textbf{Lasers (1972):} $f$ was counted with a caesium atomic clock; $\lambda$ was
        compared with krypton light. An indirect method: $c = f\lambda$.
  \item \textbf{A school method:} in a microwave oven with no turntable, the microwaves make a
        \term{standing wave} (a wave pattern that stays in one place). Chocolate melts at its
        \term{antinodes}: fixed spots where the wave is always strong. Melted spots are half a
        wavelength apart: $\lambda = 2\times\text{(spacing)}$. With the oven's
        $f = \qty{2.45}{GHz}$, $c = f\lambda \approx \qty{3.0e8}{m.s^{-1}}$.
  \item \textbf{Since 1983}, $c$ is \emph{defined} as exactly \qty{299792458}{m.s^{-1}} (the
        krypton metre had limited the laser result).
        \textbf{One metre} is the distance light travels in a vacuum in $1/299\,792\,458$ of a
        second. \textbf{The second} (since 1967) is \num{9192631770} oscillations of the
        radiation from one energy jump in a caesium-133 atom.
  \item This metre (before: a metal bar, then krypton light) can be made in any laboratory and
        rests on a constant of nature. We measure lengths by timing light (radar, GPS).
\end{itemize}



%% ==================================================================
\section{Spectra and identifying elements}\label{ch1:sec:spectra}
\covers{M7.IQ1.04, M7.IQ1.05}

\begin{bigidea}
Hot, dense objects give a spectrum with all colours. Low-density gases give lines, unique to
each element. So spectra identify elements.
\end{bigidea}

\begin{keywords}
  \kw{Spectrum (plural spectra)}{light spread out in order of wavelength.}
  \kw{Spectroscope}{an instrument that spreads light into a spectrum.}
  \kw{Discharge tube}{a glass tube of gas at low pressure that glows when a high voltage is
    applied.}
  \kw{Incandescent}{glowing because it is hot.}
\end{keywords}

A spectroscope uses a prism or a \term{diffraction grating} (thousands of tiny, equal slits).
Using spectra to learn about a source is \term{spectroscopy}. \Cref{ch1:fig:spectra} shows the
three kinds (Kirchhoff's rules, 1859).

\begin{sidebyside}[0.675]
  \centering
  \includegraphics{ch1_spectra.pdf}
  \captionof{figure}{The three kinds of spectra.}\label{ch1:fig:spectra}
\nextside
  \centering
  \begin{tikzpicture}[x=0.8cm, y=1cm]
    \draw[->, axisline] (-0.55,0) -- (-0.55,2.7) node[above, font=\footnotesize] {energy};
    \foreach \y/\n in {0/1, 1.45/2, 2.1/3, 2.45/4}
      \draw[line width=0.9pt, hscink] (0,\y) -- (3.15,\y) node[right, font=\footnotesize] {$n=\n$};
    \draw[->, line width=1.1pt, hscforce] (0.95,2.1) -- (0.95,1.48);
    \draw[->, decorate, decoration={snake, amplitude=1.3pt, segment length=5pt, post length=2mm},
      line width=0.8pt, hscray] (0.85,1.78) -- (0.05,1.78);
    \node[font=\footnotesize, align=center, anchor=north] at (0.5,1.38) {emission:\\photon out};
    \draw[->, line width=1.1pt, hscvel] (2.25,1.45) -- (2.25,2.07);
    \draw[->, decorate, decoration={snake, amplitude=1.3pt, segment length=5pt, post length=2mm},
      line width=0.8pt, hscray] (3.1,1.78) -- (2.35,1.78);
    \node[font=\footnotesize, align=center, anchor=north] at (2.68,1.38) {absorption:\\photon in};
    \node[font=\footnotesize, anchor=north] at (1.6,-0.05) {lowest level};
  \end{tikzpicture}
  \captionof{figure}{Electrons jump between energy levels. Wavy arrows: photons.}%
  \label{ch1:fig:levels}
\end{sidebyside}

\begin{enumerate}
  \item \term{Continuous}: all wavelengths, no gaps. From a hot, dense object: an incandescent
        filament, lava, the inside of a star.
  \item \term{Emission}: bright lines on a dark background. From a \term{low-density} gas (atoms
        far apart) whose atoms are \term{excited} (given extra energy). A flame or a discharge
        tube excites them. Seen from any side.
  \item \term{Absorption}: a continuous spectrum with dark lines, when light from a hot dense
        source passes through a cooler, low-density gas. Seen only looking through the gas at
        the source. The dark lines are at \emph{exactly the same wavelengths} as that gas's
        bright lines.
\end{enumerate}

\begin{itemize}
  \item \textbf{Where lines come from:} electrons in an atom have only certain energies: \term{energy
        levels} (\cref{ch1:fig:levels}). An electron that drops a level gives out a
        \term{photon} (a packet of light energy) equal to the gap. One gap gives one frequency
        ($E = hf$): one line. An atom absorbs only photons that match a gap: dark lines.
        Each element has its own levels, so its own lines; emission and absorption use the same
        gaps, so the same wavelengths.
  \item \textbf{Identifying elements:} compare the lines with laboratory lines (for example
        hydrogen: 410, 434, 486, \qty{656}{nm}). A line given as a frequency: use $c = f\lambda$.
        An element is present only if \emph{all} its strong lines are there.
  \item Helium was found in the Sun's spectrum (1868) before it was found on Earth.
\end{itemize}

%% ==================================================================
\section{What starlight tells us}\label{ch1:sec:stars}
\covers{M7.IQ1.06}

\begin{bigidea}
The peak of a star's spectrum gives its temperature; its lines give its composition. Shifted
lines show motion; wider lines show rotation or density.
\end{bigidea}

\begin{keywords}
  \kw{Black body}{an ideal object that absorbs all radiation and is a perfect emitter. Its
    spectrum depends only on its temperature.}
  \kw{Doppler effect}{the change in wavelength when a source and an observer move towards or
    away from each other.}
  \kw{Redshift, blueshift}{lines moved to longer (red) or shorter (blue) wavelengths.}
  \kw{Translational velocity}{how fast the whole star moves through space.}
  \kw{Rotational velocity}{how fast the star spins (the speed of its surface as it turns).}
  \kw{Line broadening}{spectral lines becoming wider.}
\end{keywords}

A star's hot, dense inside gives a continuous spectrum. Atoms in its cooler, lower-density
outer layers (its \term{atmosphere}) absorb the wavelengths that match their energy gaps. So a
star gives an absorption spectrum (in sunlight: the \term{Fraunhofer lines}). Its continuous
part is close to the spectrum of a black body.

\textbf{1. Chemical composition:} match the dark lines with laboratory lines. The Sun's lines
show hydrogen, sodium, magnesium, calcium, iron and many more elements.

\begin{sidebyside}[0.675]
  \centering
  \includegraphics{ch1_blackbody.pdf}
  \captionof{figure}{Black-body spectra at three temperatures.}\label{ch1:fig:blackbody}
\nextside
  \textbf{2. Surface temperature} (\cref{ch1:fig:blackbody}). A hotter star:
  \begin{itemize}
    \item peaks at a shorter wavelength (Wien's law);
    \item gives out more light at every wavelength from each square metre;
    \item looks bluer.
  \end{itemize}
\end{sidebyside}

Line strengths also depend on temperature: hydrogen lines are strongest near \qty{10000}{K};
cool stars (below about \qty{3700}{K}) show wide bands made by molecules. Classes O, B, A, F, G,
K, M go from hottest to coolest.

\begin{formulacard}{Wien's displacement law}
  \cardformula{\lambda_\mathrm{max} = \frac{b}{T}}
  \begin{symbols}
    \sym{$\lambda_\mathrm{max}$}{peak wavelength of the continuous spectrum}{\si{m}}
    \sym{$b$}{Wien's displacement constant, \num{2.898e-3} (data sheet)}{\si{m.K}}
    \sym{$T$}{surface temperature}{\si{K}}
  \end{symbols}
  \whentouse{the peak of the continuous (black-body) spectrum, never one line.}
  \trap{$T$ in kelvin: $T(\unit{K}) = T(\unit{\degreeCelsius}) + 273$. Change \si{nm} to \si{m}.}
\end{formulacard}

\textbf{3. Translational velocity.} A siren sounds higher as it comes towards you and lower
as it goes away: the Doppler effect. Light does the same (\cref{ch1:fig:starlines}).
\begin{itemize}
  \item Moving away: redshift (every line moves to a longer wavelength). Moving towards us:
        blueshift. Faster star, bigger shift. The whole pattern moves, so the elements can
        still be recognised. Only motion along our \term{line of sight} shows up.
  \item $\Delta\lambda = \dfrac{\lambda v}{c}$ \hscbadgenotonsheet\ (for $v$ much smaller than
        $c$). In one period $T = \lambda/c$ (here $T$ is a period, not a temperature), a star
        moving away travels $vT$. So its crests are $\lambda + vT$ apart:
        $\Delta\lambda = vT = \lambda v/c$.
  \item Shifts are tiny, because $v$ is far below $c$. So spectroscopes must be very precise.
        The sheet's Doppler formula is for sound.
\end{itemize}

\begin{sidebyside}[0.71]
  \centering
  \includegraphics{ch1_starlines.pdf}
  \captionof{figure}{Hydrogen absorption lines in four cases.}\label{ch1:fig:starlines}
\nextside
  \textbf{4. Rotational velocity.} One edge turns towards us (blueshift), the other away
  (redshift). We see both at once, so each line is broadened, but its centre does not move.
  Faster spin, wider lines.
\end{sidebyside}

\textbf{5. Density.} In a dense atmosphere, atoms collide often and sit close to other charged
particles. This disturbs their energy levels, so the lines are wider (\term{pressure
broadening}). \term{Giant stars} (huge, swollen, near the end of their lives) have low-density
outer layers and narrow lines; \term{main-sequence stars} such as the Sun are denser.

%% ==================================================================
\section{Key experiments}\label{ch1:sec:experiments}
\covers{M7.IQ1.03, M7.IQ1.04}

\begin{experiment}{Fizeau's toothed wheel (1849)}
\aim{To measure the speed of light directly, on Earth.}
\method{Light from a \term{half-silvered mirror} (it reflects half the light and lets half
  through) passes a gap in a \term{toothed wheel} ($N = 720$ teeth and 720 gaps). It goes to a
  mirror \qty{8.63}{km} away and back (\cref{ch1:fig:fizeau}). The wheel spins faster until the
  returning light first hits the next tooth: darkness.
  \begin{figure}
    \centering
    \begin{tikzpicture}[x=1cm, y=1cm]
      \draw[line width=0.8pt] (0,0) circle (0.2);
      \fill[hscink] (0.09,0) circle (0.07);
      \node[font=\footnotesize, anchor=north] at (0,-0.25) {observer};
      \draw[line width=1.6pt, hscfield!80] (1.1,-0.28) -- (1.66,0.28);
      \node[font=\footnotesize, align=center, anchor=south] at (1.38,0.35) {half-silvered\\mirror};
      \draw[line width=0.8pt, fill=hscray!35] (1.38,-1.3) circle (0.2);
      \node[font=\footnotesize, anchor=east] at (1.12,-1.3) {lamp};
      \draw[ray] (1.38,-1.07) -- (1.38,-0.07);
      \foreach \y in {-0.95,-0.55,0.25,0.65}
        \fill[hscink] (2.9,\y) rectangle (3.08,\y+0.3);
      \draw[line width=0.6pt, hscmuted] (2.99,-1.05) -- (2.99,1.05);
      \node[font=\footnotesize, align=left, anchor=west] at (3.15,-0.85) {wheel\\(edge-on)};
      \draw[ray] (1.44,0.06) -- (4.5,0.06);
      \draw[ray] (5.7,0.06) -- (8.1,0.06);
      \draw[ray] (8.1,-0.08) -- (5.7,-0.08);
      \draw[ray] (4.5,-0.08) -- (0.28,-0.08);
      \draw[line width=0.6pt] (4.85,-0.28) -- (5.05,0.28);
      \draw[line width=0.6pt] (5.15,-0.28) -- (5.35,0.28);
      \draw[line width=2pt, hscink] (8.2,-0.55) -- (8.2,0.55);
      \fill[ground] (8.22,-0.55) rectangle (8.42,0.55);
      \node[font=\footnotesize, anchor=north] at (8.25,-0.6) {mirror};
      \draw[<->, line width=0.6pt, hscmuted] (2.99,-1.6) -- (8.2,-1.6)
        node[midway, fill=tcbcolback, font=\footnotesize, text=hscmuted] {$D = \qty{8.63}{km}$};
      \begin{scope}[shift={(6.1,1.15)}]
        \foreach \i in {0,...,15}
          \fill[hscink] ({22.5*\i-5.6}:0.45) -- ({22.5*\i-5.6}:0.62)
            arc ({22.5*\i-5.6}:{22.5*\i+5.6}:0.62) -- ({22.5*\i+5.6}:0.45)
            arc ({22.5*\i+5.6}:{22.5*\i-5.6}:0.45) -- cycle;
        \draw[line width=0.8pt, fill=hscmuted!20] (0,0) circle (0.45);
        \fill[hscink] (0,0) circle (0.05);
        \node[font=\footnotesize, align=left, anchor=west] at (0.72,0) {wheel face-on:\\$N$ teeth, $N$ gaps};
      \end{scope}
    \end{tikzpicture}
    \caption{Fizeau's apparatus. Light leaves through a gap. At the right speed, it returns to
      find a tooth.}\label{ch1:fig:fizeau}
  \end{figure}}
\results{Darkness at $f \approx 12.6$ turns per second (here $f$ is how fast the wheel turns,
  not a light frequency). In the light's travel time, a gap turns to the next tooth: $1/(2N)$ of
  a turn. One turn takes $1/f$. \textbf{Travel time} \hscbadgenotonsheet\ $t = 1/(2Nf)$. With
  $d = 2D$: $c \approx \qty{3.13e8}{m.s^{-1}}$.}
\conclusion{$c \approx \qty{3.1e8}{m.s^{-1}}$, within about 5\,\% of today's value.}
\significance{The first measurement of $c$ on Earth. Maxwell compared his prediction with
  values like this.}
\end{experiment}

\begin{experiment}{Spectra of discharge tubes, a lamp and the Sun}
\aim{To observe and compare the spectra of different light sources.}
\method{In a dark room, view H, He, Ne and Hg discharge tubes through a spectroscope. Then view
  a filament lamp, and sunlight reflected from a white card. Read each line's wavelength.
  Safety: thousands of volts; mercury tubes give out UV; never look at the Sun.}
\results{{\centering\small
  \begin{tabular}{@{}L{32mm}L{98mm}@{}}
    \toprule
    \colhead{Source} & \colhead{Spectrum seen}\\
    \midrule
    discharge tubes & emission: bright lines on black, a different set for each gas\\
    filament lamp & continuous: all colours, brightest at the red end\\
    reflected sunlight & continuous, with narrow dark lines, e.g.\ \qty{589}{nm} (Na), \qty{656}{nm} (H)\\
    \bottomrule
  \end{tabular}\par}}
\conclusion{Hot dense sources give continuous spectra. Excited gases give line spectra,
  unique to each element.}
\significance{The basis of spectroscopy in the laboratory and in astronomy.}
\end{experiment}

\begin{summarybox}{Remember}
\begin{itemize}
  \item Maxwell: a changing electric field makes a magnetic field (unification). He predicted EM
        waves at $1/\sqrt{\mu_0\varepsilon_0} = c$, so light is an EM wave. Hertz confirmed it.
  \item Accelerating charges make EM waves: $E$ and $B$ in phase, at \ang{90} to each other and
        to the direction of travel, no medium. All travel at $c = f\lambda$.
  \item Since 1983, $c = \qty{299792458}{m.s^{-1}}$ exactly; the metre comes from $c$ and the
        caesium second.
  \item Continuous: hot dense source. Emission: excited low-density gas. Absorption: cool gas in
        front of a hot source. Same element, same wavelengths.
  \item Star: peak $\to$ temperature ($T$ in kelvin); lines $\to$ composition; shift $\to$ motion
        along the line of sight; broadening $\to$ rotation or density.
\end{itemize}
\end{summarybox}
